\section{Framework}
\label{sec:method}

A recursive self-improvement loop for a robot policy makes three decisions every round (\cref{fig:framework}): which policy generates the experience it trains on, which trained candidate it keeps, and whether it revises the procedure that trains it. The kept version seeds the next round: its weights generate the next experience and its program trains the next candidates. The loop thus improves two layers, the \emph{artifact} (the policy that acts) and the \emph{harness} (the training program that produces it). \Cref{sec:experience,sec:selection,sec:meta} define the three decisions; \cref{sec:controls} turns each into a switch.

\paragraph{Setting.} A task instance $c=(\ell, x_0)$ is a language instruction $\ell$ and an initial configuration $x_0$ of the scene (the object placements; we also call it a \emph{scene}). The policy $\pi_\theta$ maps the instruction and its camera and proprioceptive observations to actions. Rolling it out in the simulator from $x_0$ produces a trajectory $\tau$, whose distribution depends on the policy, the environment dynamics and the sampling temperature. Each trajectory is labelled by the task's fixed success rule $v(c, \tau) \in \{0, 1\}$, which neither the policy nor the program proposer can modify. A \emph{training program} $P$ specifies how experience is collected (which instances, $K$ attempts per instance, temperature $T$), which attempts become training data, and how the policy is updated; it is stored as a JSON configuration, and $\mathcal{P}$ is a declared space of allowed values. A \emph{version} is a pair $(\theta, P)$; the version the loop currently keeps is the \emph{retained} version $(\theta^\star, P^\star)$.

\subsection{Learning from self-generated experience}
\label{sec:experience}
A \emph{training block} runs program $P$ on policy $\theta$: it samples task instances, rolls out $K$ attempts on each, obtains their success labels, turns the attempts into training data, and updates a LoRA adapter on $\theta$ with the configured optimiser steps. Two data rules are used: \texttt{mixed\_groups} keeps every attempt of instances whose success rate lies strictly inside a band, for GRPO~\cite{shao2024deepseekmath} (the group-centred advantage is zero when all attempts agree), and \texttt{verified} keeps successful attempts only, for behaviour cloning (BC). A \emph{round} trains each candidate for three successive blocks: every candidate starts from the same retained weights $\theta^\star$, and within a candidate each block continues from the weights of the previous block (only the adapter carries over; each block starts a fresh optimiser). In the recursive configuration each block collects experience with the current candidate (the first with $\theta^\star$), so experience always comes from the latest policy: the \emph{data recursion}. The control \texttt{nonrec} collects every block's experience with the frozen initial policy $\theta_0$ while still updating the learner from $\theta^\star$.

\begin{algorithm}[t]
\caption{Gated recursive self-improvement}
\label{alg:loop}
\begin{algorithmic}[1]
\REQUIRE $\theta_0$, $P_0$, success rule $v$, rounds $R$, candidates $c$
\STATE $(\theta^\star, P^\star) \leftarrow (\theta_0, P_0)$
\FOR{$r = 1 \ldots R$}
  \STATE $\{P_i\}_{i \le c} \leftarrow \{P^\star\} \cup \textsc{Propose}(\text{records}_{<r})$ \hfill $\triangleright$ meta
  \STATE draw fresh $G_r$; $g^\star \leftarrow$ success of $\theta^\star$ on $G_r$
  \FOR{each $P_i$}
     \STATE $\theta_i \leftarrow \theta^\star$
     \FOR{block $= 1 \ldots 3$}
        \STATE $\theta_{\text{gen}} \leftarrow \theta_i$ (recursive) or $\theta_0$ (\texttt{nonrec})
        \STATE attempts $\leftarrow$ $K$ rollouts of $\theta_{\text{gen}}$ on \textsc{Sample}$(P_i)$
        \STATE labels $\leftarrow v$ (\texttt{shuffled}: permuted per instance)
        \STATE $\theta_i \leftarrow \textsc{Update}(\theta_i, \textsc{Select}(\text{attempts}, P_i), P_i)$
     \ENDFOR
     \STATE $g_i \leftarrow$ success of $\theta_i$ on $G_r$ (greedy)
  \ENDFOR
  \STATE $i^\star \leftarrow \arg\max_i g_i$; keep $(\theta_{i^\star}, P_{i^\star})$ if $g_{i^\star} \ge g^\star$
\ENDFOR
\STATE \textbf{return} $(\theta^\star, P^\star)$ \hfill $\triangleright$ tests read after all runs end
\end{algorithmic}
\end{algorithm}

\subsection{Selecting policy updates}
\label{sec:selection}
The loop compares each trained candidate with its parent on the same newly sampled validation instances, the \emph{gate}: 20 instances per task, 200 episodes, greedy decoding, the environment re-seeded per episode so that both versions face identical initial configurations. The best candidate and its program are retained if its success rate is at least the parent's; otherwise the parent version is kept. A new validation set is drawn every round and the parent is re-evaluated on it, so every decision is a paired comparison on instances that no version has trained or been selected on. This avoids repeated selection against one fixed set, but each decision remains subject to finite-sample noise (\cref{sec:rq2}). The control \texttt{nogate} accepts every candidate. An audit set and two sealed tests play no part in selection (\cref{sec:protocol}).

\subsection{Adapting training programs}
\label{sec:meta}
The \emph{meta level} revises the harness. At the start of each round a proposer reads the records of completed rounds and returns $c{-}1$ programs inside $\mathcal{P}$; it does not see the current round's validation scores. The incumbent program $P^\star$ trains a continuation candidate and each proposed program an alternative one, all from the same parent weights, and the gate selects among them; a program is retained only together with the candidate it trained. With the incumbent alone, adaptation is off (\texttt{fixed}); otherwise one of three proposers writes the programs:
\begin{itemize}[nosep,leftmargin=*]
  \item \texttt{full} (numeric search): random mutation of declared knobs, biased toward knobs whose past changes raised the gate score.
  \item \texttt{llm}: a language model (Claude Sonnet) reads the current program, the edit space with a neutral description of what each option computes, and, for every past candidate, its program diff, gate score, acceptance and per-block statistics (sampled success, share of mixed instances, data counts, first and last training loss). It writes a diagnosis, then returns one program.
  \item \texttt{llm\_harness}: the same model with a proposer workspace, following~\cite{lee2026metaharness,chen2026evotrainer}: files with the full record of every candidate, per-task gate tables, an evidence ledger pooling gate changes per edited field, a frozen memory of earlier runs with a playbook distilled before launch, and its own notes, read with list, read and search tools over several turns; the prompt states the measured gate noise. The workspace informs the proposal; what is edited is still the training program. Building it exposed a field-name bug: the plain \texttt{llm} proposer never received the training loss it was told to watch.
\end{itemize}
The \emph{knob} space has ten hyperparameters (instances, allocation, samples, temperature, data rule and budget, operator, learning rate, steps, batch); the \emph{harness} space adds nine components of the procedure (policy-gradient ratio, loss aggregation, BC likelihood temperature, advantage normalisation, data band, per-instance cap, allocation floor, two clip bounds). Each is a change made by hand during development after diagnosing round records; the LLM is never told which option is better.

\paragraph{Attribution guard.} The proposer never sees the audit, the tests, the weights or the episodes; it only chooses values inside $\mathcal{P}$, every proposal passes the same validator as a numeric mutation, failures fall back to numeric search and are counted, and an optional \emph{compute cap} keeps any candidate within the seed program's rollout and optimiser budget. All training data comes from the policy's own attempts and every label from the success rule; the LLM chooses the procedure, not the data.
